\documentclass[11pt]{article}
\usepackage[T1]{fontenc}
\usepackage{amsmath,amssymb,amsthm}
\usepackage[margin=1in]{geometry}
\usepackage[hidelinks]{hyperref}
\newtheorem{theorem}{Theorem}
\newtheorem{lemma}[theorem]{Lemma}
\newtheorem{proposition}[theorem]{Proposition}
\newcommand{\DD}{d'd''}
\newcommand{\MA}{\operatorname{MA}}
\newcommand{\C}{\mathcal C}
\title{Angular degeneration of nonskeletal Dirac equations\\
An isotropic dual bound and failure of unordered pointwise continuity}
\author{Independent scoped P16.7 deductions; unreviewed working paper}
\date{9 October 2026}
\begin{document}
\maketitle
\begin{abstract}
Varying the residue angle of a quadratic pencil produces explicit
continuous positive weights with Monge--Amp\`ere measures twice an
anisotropic Gauss point. Those measures converge to twice the original
isotropic Gauss point. Weak lower semicontinuity of the full dual then
improves the previously retained upper bound: at depth $1/8$, the dual
lies between $3/8$ and $1/2$. The potentials converge pointwise to a
different one-well weight, whose measure and energy differ from their
limits. Their arithmetic means converge uniformly, with exact error
$a/N$. Finally, the individual pencils have a dual objective gap at
least $a$ for the limiting target. Thus weak approximation by solved
equations supplies no maximizing sequence for that target. All claims
retain the disclosed bounded-calculus and geometric qualifications;
existence for the isotropic whole Dirac and the original problem remain open.
\end{abstract}

\section{Exact datum and retained analytic boundary}
Keep the supplied compact, good, strict, smooth, separated, boundaryless
pure surface
\[
 K=\mathbb C((s)),\quad |s|=e^{-1/8},\quad t=s^8,\qquad
 X=(\mathbb G_m^2)^{\rm an}/\langle(t,2),(3,t)\rangle.
\]
The geometry is an input. Genuine nonalgebraization and coherent descent
are not independently certified. On an injective ordinary chart
$U=\{|r_i|<1/4\}$, let $r_i=-\log|z_i|$, $\rho=|r|^2/2$,
$\theta=\DD\rho$, and $g=\sigma(0)$ on the canonical Gauss section.
The class $\C$ is the entire original minimal continuous CLD class of
weights with $\rho+q$ locally psh. The raw mass is two, and
\[
 E(q)=\frac13\int_Xq(\theta^2+\theta\theta_q+\theta_q^2),\qquad
 D(\mu)=\sup_{q\in\C}\left(E(q)-\int_Xq\,d\mu\right).
\]
The generated class requires actual decreasing continuous presentations;
an unordered pointwise limit is not such an antecedent.

The qualified bounded calculus is reconstructed in the retained whole
ROOT436 source: positive recursive products, smooth-gradient closability
and local uniform closure for the entire minimal sheaf, global Stokes,
constant mixed mass two, and cubic first variation and concavity. In
particular, for every actual bounded $q\in\C$,
\[
 E(p)-\int p\,\MA(q)\le E(q)-\int q\,\MA(q)\quad(p\in\C).       \tag{1}
\]
The baseline and coordinate-fibre computation are also reconstructed there.
These are written deductions, not incoming correctness grants. Acceptance,
the whole historical audit, deeper analytic foundations and genuine
geometry remain unverified. The complete ROOT434 and ROOT436 sources and
their qualifications accompany this paper.

The essential external results are from Chambert--Loir--Ducros,
\emph{Formes diff\'erentielles r\'eelles et courants sur les espaces de
Berkovich}, version 2, 25 July 2025. Definition 8.1.7, Lemma 8.1.8,
Corollary 8.2.4, Sections 8.5, 8.7 and 8.8 supply actual minimal-class
positivity, convex maxima, recursive positive products, local uniform
continuity, proper fixed-fibre restriction and multilinearity. Proposition
5.11.1 supplies smooth Stokes on the stated empty-boundary compact domain;
the whole-class extension is the retained gradient argument. Sections
7.2.8--13 supply Radon reconstruction and all-continuous-test equality.
The complete immediate 6.2.1--9 chain was additionally read for ROOT436;
smooth cutoffs, density and partitions apply on our locally differentially
separated ordinary opens and paracompact compact Hausdorff carrier.
Appendix A and deeper recursive foundations keep their qualification.

Corollary 18.6.2 identifies
$(\DD\log\max(R,|f_1|,|f_2|))^2$ with the integration current of the
whole fibre over the radius-$(R,R)$ Gauss point, for ordinary boundaryless,
topologically separated, equidimensional analytic spaces and arbitrary
positive real $R$. It does not assign probability mass one to each point.
Its complete Proposition 18.6.1 proof and the finite-flat domain in
11.6.1--5 were read and checked for ROOT434. Below the map is explicitly
finite flat of degree two; translations and unit linear changes have
degree one. No value-group restriction or primitive degree-one replacement
is used. Literal prose slips about first-current vanishing are avoided:
the rank argument only gives vanishing of the top power.

\section{Pencils at every residue angle}
Put $u=z_1-1$, $v=z_2-1$, $W=\{|u|<1,|v|<1\}$ and
\[
 q_0=\sum_{i=1}^2 \{r_i\}(1-\{r_i\})/2-1/4,
 \quad A_i=\max(1,|z_i|),\quad H_0=\log(A_1A_2).
\]
Then $q_0\in\C$, $\MA(q_0)=2\delta_g$, $E(q_0)=-1/3$, and
$\rho+q_0=H_0-\frac12\log|z_1z_2|-1/4$ on $U$.
For any constant $\alpha\in\mathbb C^*$ and depth $c>0$, set
\[
 H_{c,\alpha}=\log\max(|uv|,|u+\alpha v|,e^{-c}),\qquad
 Q_{c,\alpha}=q_0+H_{c,\alpha}-H_0
\]
on $U$, extending the correction by zero outside $U$.
Let $\xi_{c,\alpha}$ be the joint Gauss point of radii
$(e^{-c},e^{-c/2})$ in
\[
 x=u+\alpha v,\qquad y=u-\alpha v.
\]

\begin{proposition}\label{pencil}
For every $c>0$ and $\alpha\in\mathbb C^*$, the weight is an actual
bounded member of the original $\C$, and
\[
 \MA(Q_{c,\alpha})=2\delta_{\xi_{c,\alpha}},\quad
 E(Q_{c,\alpha})=-1/3-2c/3,\quad
 D(2\delta_{\xi_{c,\alpha}})=1/6+4c/3.
\]
\end{proposition}
\begin{proof}
All constants $2,4,\alpha$ have norm one. Outside $W$, if both translated
norms exceed one, their product dominates the sum. If exactly one exceeds
one and the other is below one, the sum has the larger norm; if the other
equals one the product has it. If both translated norms are at most one
and at least one equals one, their product has norm one when both do,
and their sum does when only one does. These cases also give the exact
original-coordinate value $A_1A_2$ when a $|z_i|$ is below one, since
then its translated norm is one. Thus
$\max(|uv|,|u+\alpha v|)=A_1A_2$ outside $W$, on every actual seminorm.

The positive constant slot permits omission of any zero polynomial slot
on a neighborhood. The remaining unit logarithms admit smooth convex
positive maxima converging locally uniformly. Subtracting the same smooth
harmonic reference term proves actual positivity. The correction is zero
on the ordinary collar $1/8<\max|r_i|<1/4$ and glues to $q_0$.
The deck units $2,3$ are retained; no translated-coordinate deck invariance
is presumed. On $W$, $4\alpha uv=x^2-y^2$ and $|x|<1$, giving
\[
 H_{c,\alpha}=\log\max(e^{-c},|x|,|y|^2).
\]
The determinant $-2\alpha$ is a unit. The actual map
$(x,y)\mapsto(x,y^2)$ between ordinary unit bidiscs is finite flat,
with module basis $1,y$ and degree two. Over the target Gauss completed
residue field $H$, its fibre algebra is $H[T]/(T^2-Y)$. In the graded
Gauss fraction field, the independent variable $Y$ has odd divisor order,
so is not a square. Completion preserves initial forms by strictly closer
approximation. Hence this algebra is a field, with unique extended
absolute value. The Gauss norm of radii $(e^{-c},e^{-c/2})$ restricts to
the target norm, so the entire fibre consists of $\xi_{c,\alpha}$.

Corollary 18.6.2 puts all the local top measure on this point. Outside
$W$ the cap is inactive; at zero slots omit them, and with two units the
maximum is a smooth null log plus a one-ratio convex maximum, with zero
top power. This covers $g$; away from the zero tropical fibre the weight
agrees locally with $q_0$. The collar contributes no measure. Positivity
and global Stokes total mass two therefore give exactly coefficient two.

The fixed original-coordinate unit Gauss curves restrict this maximum
to the opposite baseline coordinate maximum: on the first curve
$|u|=1$ and $\max(|uv|,|u+\alpha v|)=\max(1,|v|)=A_2$, and conversely
on the second. Proper fixed-fibre restriction gives
$B(q_0,Q_{c,\alpha})=2\delta_g$. Thus, writing $w=Q_{c,\alpha}-q_0$,
\[
 \MA(q_0+h w)=2(1-h^2)\delta_g+2h^2\delta_{\xi_{c,\alpha}}.
\]
Here $w(g)=0$, $w(\xi)=-c$, and $q_0=-1/4$ at both. First variation
gives $E=-1/3-2ch^3/3$; the target moment at $h=1$ is $-1/2-2c$.
Inequality (1) on the entire original class computes its full dual.
\end{proof}

\section{Weak degeneration and the isotropic full-dual bound}
Fix $a>0$, let $r=e^{-a}$, choose distinct constants
$\alpha_j=j+1$, and abbreviate
\[
 \xi_j=\xi_{2a,\alpha_j},\quad Q_j=Q_{2a,\alpha_j},\quad
 \mu_j=2\delta_{\xi_j}.
\]
Let $x_a$ denote the ORIGINAL isotropic joint Gauss point of radii $(r,r)$
in $(u,v)$, and $\mu=2\delta_{x_a}$.

\begin{lemma}\label{weak}
In the actual ordinary analytic topology, $\xi_j\to x_a$, hence
$\mu_j\to\mu$ weakly on every continuous analytic test.
\end{lemma}
\begin{proof}
For a nonzero polynomial $P(u,v)=\sum c_{ik}u^iv^k$, its isotropic norm is
$M=\max |c_{ik}|r^{i+k}$. Substitution
$u=(x+y)/2$, $v=(x-y)/(2\alpha)$ has norm-one coefficients and
$|x|=r^2\le r$, $|y|=r$, so its anisotropic norm is at most $M$.
Choose a degree $d$ with $m_dr^d=M$, where
$m_d=\max_{i+k=d}|c_{ik}|$. Choose a coefficient of norm $m_d$ to
normalize this homogeneous part. The normalized coefficient of $y^d$
after substitution has residue
\[
 2^{-d}\sum_{i+k=d}\overline{c_{ik}/c_*}(-\alpha^{-1})^k.
\]
This is a nonzero Laurent polynomial in the residue of $\alpha$:
the exponents $k$ are distinct and some coefficient is nonzero. It
vanishes at only finitely many of our distinct residues. Otherwise the
coefficient has norm exactly $m_d$, and its $y^d$ monomial has norm $M$.
Other homogeneous degrees cannot change that coefficient. The Gauss norm
is therefore exactly $M$ for all but finitely many $j$.

Both original coordinates have norm one at these points. Multiplying a
Laurent polynomial in $z_1,z_2$ by a monomial gives an ordinary polynomial
in $u,v$ without changing its norm. Thus each fixed regular function on
the algebraic torus has its isotropic norm eventually. The defining
seminorm topology of the torus analytification gives convergence in the
chart, then in $X$. Integration of each continuous function at a convergent
point proves the weak measure assertion. No countable determining family
or value-group assumption is used.
\end{proof}

\begin{theorem}\label{bound}
For every $a>0$,
\[
 \frac16+\frac{5a}3\ \le\ D(2\delta_{x_a})\ \le\
 \frac16+\frac{8a}3.
\]
In particular, for the originally prescribed depth $a=1/8$,
\[
                         3/8\le D(2\delta_{x_{1/8}})\le1/2.
\]
\end{theorem}
\begin{proof}
For every fixed $p\in\C$, the function
$\lambda\mapsto E(p)-\int p\,d\lambda$ is weakly continuous.
Their supremum $D$ is therefore weakly lower semicontinuous. Lemma
\ref{weak} and Proposition \ref{pencil} give
\[
 D(\mu)\le\liminf_jD(\mu_j)=1/6+8a/3.
\]
For the lower bound, define
$\gamma_{\rm iso}=\max(0,-\log\max(|u|,|v|))$ on $U$, zero outside,
and $q_b=q_0-\min(b,\gamma_{\rm iso})$. Its actual positive local maximum
is $\log\max(e^{-b},|u|,|v|,|uv|)$ minus the same harmonic reference.
The translated Gauss fibre at depth $b$ and inverse-coordinate fibre at
$g$, each degree one, together with positive total mass two give
$\MA(q_b)=\delta_g+\delta_{x_b}$ and $B(q_0,q_b)=2\delta_g$.
First variation gives $E(q_b)=-1/3-b/3$. For $b=a$, the moment against
$\mu$ is $-1/2-2a$, so this actual competitor gives $1/6+5a/3$.
This upper bound establishes finiteness, not attainment or existence.
\end{proof}

\section{Unordered pointwise limits and uniform arithmetic means}
Put $q_*=q_{2a}$, the isotropic one-well weight just defined.
\begin{theorem}\label{unordered}
The $Q_j$ are uniformly bounded actual members of $\C$ and are pairwise
incomparable. On EVERY actual point, $Q_j\to q_*$, in fact with at most
one exceptional index. Nevertheless
\[
 \MA(Q_j)\longrightarrow2\delta_{x_a}\ne
 \MA(q_*)=\delta_g+\delta_{x_{2a}},
\]
\[
 E(Q_j)=-1/3-4a/3\ne E(q_*)=-1/3-2a/3.
\]
For the actual convex averages $A_N=N^{-1}\sum_{j=1}^NQ_j$,
\[
 \|A_N-q_*\|_\infty=a/N,\quad
 \MA(A_N)\longrightarrow\delta_g+\delta_{x_{2a}},\quad
 |E(A_N)-E(q_*)|\le2a/N.
\]
\end{theorem}
\begin{proof}
For distinct $j,k$, the linear system with rows $(1,\alpha_j)$ and
$(1,\alpha_k)$ and its inverse have norm-one coefficients. Thus at every
analytic point
\[
 \max(|u+\alpha_jv|,|u+\alpha_kv|)=\max(|u|,|v|).
\]
Taking the cap and product slots gives $\max(Q_j,Q_k)=q_*$ on $U$,
and outside $U$ all equal $q_0$. Also $Q_j\le q_*$. Hence at any point
at most one $Q_j$ is strictly smaller; all others equal $q_*$. This
proves the stated pointwise convergence without monotonicity.

On $W$, if $|u|\ne|v|$, then $|u+\alpha_jv|=\max(|u|,|v|)$ and
there is no difference. If both equal $e^{-d}$, the product slot gives
\[
 0\le q_*-Q_j\le\min(2a,2d)-\min(2a,d)\le a.
\]
At zero norms use the cap directly. Outside $W$ all corrections vanish.
At $\xi_j$, the two translated norms are $r$, whereas
$|u+\alpha_jv|=r^2$. Thus $q_*-Q_j=a$ there and $Q_k=q_*$ there for
every $k\ne j$. At $\xi_k$ the inequalities reverse. This proves exact
supremum distance $a$ and pairwise incomparability, as well as uniform
bounds independent of $j$.

Proposition \ref{pencil}, Lemma \ref{weak} and the one-well calculation
give the claimed different measures and energies. The three points
$g,x_a,x_{2a}$ are distinct, so those Radon measures cannot agree.
At most one summand in $\sum(q_*-Q_j)$ is nonzero at any point; its
value is at most $a$, and is $a$ at any $\xi_j$ included in the average.
This proves the exact uniform error $a/N$. Convex positivity gives
$A_N\in\C$. Actual uniform recursive-product continuity gives its measure
limit. The cubic cocycle and positive mixed masses two bound the energy
difference by $2\|A_N-q_*\|_\infty$, proving the displayed rate.
\end{proof}

\begin{proposition}\label{gap}
The solved equations $\MA(Q_j)=\mu_j\to\mu$ do not form a maximizing
sequence for the limiting full-dual problem. Precisely,
\[
 E(Q_j)-\int Q_j\,d\mu=1/6+2a/3,\qquad
 D(\mu)-\left(E(Q_j)-\int Q_j\,d\mu\right)\ge a.
\]
\end{proposition}
\begin{proof}
At the isotropic Gauss point, $|u+\alpha_jv|=r$ and $|uv|=r^2$,
so $Q_j(x_a)=-1/4-a$. The moment is $-1/2-2a$, while its energy
is $-1/3-4a/3$. Subtract and use Theorem \ref{bound}'s lower bound.
Adding any constant changes energy and moment by the same twice-constant,
so normalization cannot remove this gap.
\end{proof}

Theorem \ref{unordered} is a failure of unordered, uniformly bounded
pointwise continuity. It does not contradict local uniform product
continuity or a theorem for actual monotone presentations. Proposition
\ref{gap} gives a concrete missing antecedent in the weak-equation
approximation route: a maximizing or sufficiently energy-controlled
sequence is still required. The specified isotropic whole-Dirac equation,
general eligible-support existence and comparison, genuine nonalgebraization
and the full original P16.7 remain open. No new acceptance or atomic
closure is claimed, and compilation remains unverified until an actual
built-in compiler success is observed.
\end{document}
